% Sustituciones por juego de los textos fijos de lang-es.tex (\renewcommand).
\renewcommand{\tComputerText}{\textbf{Cualquier MSX1} con dos ranuras de cartucho, una para el juego y otra para el V9990, y una tercera si añades un Moonsound. \tExpanderText}
\renewcommand{\tSoundText}{Se usa el PSG del MSX. Un \textbf{Moonsound} (OPL4) es opcional: añade música y voces de la recreativa. Con más de 128~KB de RAM de muestras en el Moonsound, algunos efectos también usan sonidos de la recreativa.}
\renewcommand{\tControlsText}{El teclado, para uno o dos jugadores. Los joysticks son opcionales: el puerto 1 para el jugador 1 y el puerto 2 para el jugador 2.}
\renewcommand{\tSetV}{Inserta el \textbf{cartucho V9990} en una ranura libre y, si lo tienes, el \textbf{Moonsound} en otra. \tExpanderText}
\renewcommand{\tSetGame}{Inserta el \textbf{cartucho de \gameTitle{}} en otra ranura libre. El orden no importa.}
\renewcommand{\tSetJoy}{Conecta joysticks al puerto 1 (jugador 1) y al puerto 2 (jugador 2) si quieres usarlos; el teclado sirve para los dos jugadores igualmente.}
\renewcommand{\tSetOn}{Enciende el MSX. En el monitor del V9990 aparece el logo de DiHalt Studio. Si la pantalla del propio MSX dice \texttt{This cartridge needs a V9990.} falta el V9990 o no está bien insertado, y el juego se detiene ahí.}
\renewcommand{\tPowerOnCaption}{Izquierda: la portada. Derecha: la pantalla de título.}
\renewcommand{\tCreditsTitle}{11. Créditos y derechos}
\newcommand{\tCreditsExtra}{\item \textbf{\gameTitle} fue primero una recreativa de \textbf{Atari Games}, 1987. La pantalla de título de la versión de MSX acredita \texttt{Z80 BY B. COSTAS} y \texttt{68000 BY J.TRIPP}.
\item Con un Moonsound, la música y las voces son las de la recreativa, adaptadas; la canción de la pantalla del logo es de DiHalt Studio.}
